% !TEX root = ../main.tex
\chapter{仕様・実装・テスト・証明の接続}
\label{chap:spec-impl-test-proof}

\begin{chapterabstract}
本章では、仕様、実装、テスト、証明を別々の活動としてではなく、一つの品質保証の流れとして接続する。Lean 4 による証明は、すべての品質問題を解決する万能薬ではない。外部API、UI、性能、運用環境、障害時挙動などは、通常のテスト、監視、レビューのほうが向いている。一方で、round-trip 性質、不変条件保存、件数保存、処理の意味保存のように、すべての入力について確認したい小さな核は Lean で扱いやすい。本章の目的は、何を Lean で証明し、何を通常のテストに残し、両者を設計文書の中でどう結びつけるかを判断できるようにすることである。
\end{chapterabstract}

\begin{learninggoals}
\item Lean で証明する価値が高い性質を見分けられる。
\item 通常のテスト、型、仕様、証明の役割分担を説明できる。
\item 自然言語仕様から、証明可能な小さな核を抽出できる。
\item Lean theorem を設計文書やレビュー項目に対応づける書き方を理解できる。
\end{learninggoals}

\section{この章を読む理由}

ここまでの章では、リファクタリングを等式として扱い、マイグレーションを同型や片方向互換として分類し、API adapter の一貫性を自然性として表した。これらはどれも、ソフトウェア上の変更を「何が保存されるべきか」という問いに置き換える作業であった。

しかし実務では、次のような判断がすぐに必要になる。

\begin{itemize}
\item この性質は Lean で証明すべきか。
\item この挙動は単体テストで十分か。
\item 外部APIの応答やDBの実挙動まで Lean に入れるべきか。
\item 証明した theorem を、設計文書やレビューでどう読めばよいか。
\end{itemize}

この判断を誤ると、Lean の導入は重くなる。何でも形式化しようとすると、仕様モデルが本番コードより複雑になり、変更するたびに保守できなくなる。逆に、テストだけに寄せすぎると、すべての入力に対して成り立つべき性質を代表例の確認に落としてしまう。

\begin{intuitionbox}
Lean は、テストを置き換えるものではない。Lean は、テストでは代表例しか見られない性質のうち、小さく、純粋で、壊れると困る核を検査する道具である。テストは実環境に近い振る舞いを見る。証明はモデル化した範囲の全ケースを見る。この二つは競合するのではなく、役割が違う。
\end{intuitionbox}

\FigurePlaceholder{fig-ch29-connection-map.png}{0.92}{仕様・実装・テスト・証明の接続地図}{fig:ch29-connection-map}

図\ref{fig:ch29-connection-map}では、仕様から実装へ直接進む線だけでなく、テストと証明を横に置いている。テストは「実装が期待例で動くか」を見る。証明は「モデル化した仕様がすべての場合に成り立つか」を見る。設計文書は、その対応関係と前提条件を人間が読める形で残す場所である。

\section{29.1 何を Lean で証明するべきか}

Lean で証明する価値が高いのは、実装のすべてではない。価値が高いのは、次の条件を多く満たす部分である。

\begin{itemize}
\item 入力と出力が小さな型として切り出せる。
\item 外部状態に依存しない純粋な関数として表せる。
\item 代表例ではなく、すべての入力について保証したい。
\item 仕様変更やリファクタリングで壊れると被害が大きい。
\item 証明したい性質が、等式、含意、不変条件、round-trip などとして短く書ける。
\end{itemize}

\begin{definitionbox}
本章では、\textbf{証明可能な核}という言葉を使う。証明可能な核とは、実務仕様全体のうち、Lean 上で小さな型、関数、述語、定理として表せる中心部分である。たとえば、料金計算の丸め規則すべてではなく、「この二つの実装は同じモデル関数に一致する」という等式だけを切り出す場合、その等式が証明可能な核になる。
\end{definitionbox}

典型的には、次のような性質が Lean に向いている。

\begin{longtable}{p{0.30\linewidth}p{0.61\linewidth}}
\toprule
\textbf{性質} & \textbf{Lean に向く理由} \\
\midrule
round-trip & 変換して戻すと元に戻る、という等式で書ける。スキーマ移行や encoder/decoder に向いている。 \\
不変条件保存 & 操作前に条件を満たすなら、操作後も条件を満たす、という含意で書ける。状態更新の核に向いている。 \\
件数保存 & \texttt{List.map} などの構造保存として書ける。バッチ移行の基本性質になる。 \\
意味保存 & 変更前後の関数がすべての入力で同じ結果を返す、という等式で書ける。リファクタリングに向いている。 \\
可換性 & adapter と業務ロジックの順序を入れ替えても同じ、という等式で書ける。自然性の小さな実務版である。 \\
\bottomrule
\end{longtable}

\subsection{小さな料金計算モデル}

まず、仕様と実装を分ける最小例を見る。ここでは、価格、割引、送料を自然数で扱う。現実の通貨、税率、丸め、クーポン期限、外部決済APIはまだ扱わない。これらを最初から入れると、証明の核が見えにくくなるからである。

\begin{leanbox}
次の Lean コードでは、\texttt{chargeSpec} を仕様モデル、\texttt{chargeImpl} を実装モデルとして分けている。定理 \texttt{charge\_impl\_matches\_spec} は、実装モデルが仕様モデルと同じ結果を返すことを、任意の入力 \texttt{i} について主張している。
\end{leanbox}

\begin{listing}[htbp]
\begin{minted}{lean4}
namespace Chapter29

structure PriceInput where
  base : Nat
  discount : Nat
  shipping : Nat

def chargeSpec (i : PriceInput) : Nat :=
  (i.base - i.discount) + i.shipping

def chargeImpl (i : PriceInput) : Nat :=
  let subtotal := i.base - i.discount
  subtotal + i.shipping

theorem charge_impl_matches_spec (i : PriceInput) :
    chargeImpl i = chargeSpec i := by
  rfl
\end{minted}
\caption{仕様モデルと実装モデルを分ける}
\label{lst:ch29-charge-spec-impl}
\end{listing}

ここで証明しているのは、料金システム全体の正しさではない。証明しているのは、\texttt{chargeImpl} という小さな実装モデルが、\texttt{chargeSpec} という小さな仕様モデルと一致することだけである。

\begin{softwarebox}
実務では、\texttt{chargeSpec} が仕様書中の数式や業務ルールに対応し、\texttt{chargeImpl} が本番コードから切り出した純粋な核に対応する。Lean の theorem は、「この核については、実装が仕様モデルからずれていない」というレビュー可能な証拠になる。
\end{softwarebox}

\section{29.2 何は通常のテストでよいか}

Lean は強力だが、現実のシステムのすべてを直接扱うには向いていない。通常のテストに残すべきものも多い。

\begin{longtable}{p{0.32\linewidth}p{0.59\linewidth}}
\toprule
\textbf{通常テストが向くもの} & \textbf{理由} \\
\midrule
外部APIとの通信 & 相手の実装、認証、レート制限、障害、応答遅延は Lean の小さなモデルだけでは保証できない。 \\
DBの実挙動 & インデックス、NULL、ロック、トランザクション分離レベル、実行計画は実DBで確認する必要がある。 \\
UIとブラウザ挙動 & レイアウト、アクセシビリティ、イベント、端末差は実行環境との結びつきが強い。 \\
性能とスケーラビリティ & 時間、メモリ、I/O、負荷時の挙動は測定が必要である。 \\
監視・運用・障害対応 & ログ、アラート、リトライ、フェイルオーバーは運用環境を含むテストが必要である。 \\
\bottomrule
\end{longtable}

通常テストは、具体例を見られる。Lean の証明は、モデル化した範囲を全称的に見られる。どちらが上かではなく、見ている対象が違う。

\begin{warningbox}
「Lean で証明したからテストは不要」と言ってはいけない。Lean で証明したのは、Lean に書いたモデルと定理である。本番実装がそのモデルと同期していること、外部サービスが契約通り動くこと、性能要件を満たすことは、別の確認が必要である。
\end{warningbox}

先ほどの料金計算でも、テストは役に立つ。たとえば、代表的な入力について期待値を確認することは、仕様レビューや回帰テストとして意味がある。

\begin{listing}[htbp]
\begin{minted}{lean4}
#eval chargeImpl { base := 100, discount := 10, shipping := 5 }  -- => 95

example :
    chargeImpl { base := 100, discount := 10, shipping := 5 } = 95 := by
  rfl

example :
    chargeImpl { base := 8, discount := 10, shipping := 5 } = 5 := by
  rfl
\end{minted}
\caption{代表例を確認するテスト的な使い方}
\label{lst:ch29-charge-examples}
\end{listing}

二つ目の例では、\texttt{Nat} の引き算がゼロで止まるため、\texttt{8 - 10} は \texttt{0} として扱われる。これは現実の料金計算では注意が必要な点である。このような例は、Lean の theorem とは別に、モデルの前提を人間に気づかせるテストケースとして有用である。

\begin{intuitionbox}
テストは「この入力ではどうなるか」を明らかにする。証明は「この形でモデル化したすべての入力で何が成り立つか」を明らかにする。サンプルテストは、証明の代替ではないが、モデルの読み方を確認する入口になる。
\end{intuitionbox}

\section{29.3 仕様の過剰形式化を避ける}

Lean 導入でよくある失敗は、最初から仕様全体を形式化しようとすることである。自然言語仕様には、外部依存、運用判断、業務上の例外、将来変更されるルールが含まれる。これらをすべて Lean に入れると、証明作業は仕様整理ではなく、巨大な再実装になる。

\begin{definitionbox}
\textbf{過剰形式化}とは、証明する価値に比べて、モデル化・証明・保守のコストが大きすぎる状態である。特に、変更頻度が高い周辺仕様、外部システム依存、UI表示、性能要件を、最初からすべて Lean に入れようとすると起きやすい。
\end{definitionbox}

過剰形式化を避けるには、次の三つを分ける。

\begin{enumerate}
\item \textbf{証明する核}：Lean theorem として残す部分。
\item \textbf{テストする周辺}：実装・外部環境・性能・UIなどを確認する部分。
\item \textbf{文書化する前提}：Lean モデルが何を抽象化し、何を捨てたか。
\end{enumerate}

\FigurePlaceholder{fig-ch29-proof-kernel.png}{0.90}{実務仕様から証明可能な核を切り出す}{fig:ch29-proof-kernel}

図\ref{fig:ch29-proof-kernel}のように、実務仕様全体をそのまま Lean に入れるのではなく、中心にある小さな性質を抜き出す。周辺にある環境依存の振る舞いは、テストや監視に残す。

\begin{softwarebox}
たとえば「注文金額を正しく計算する」という仕様全体には、商品価格、割引、税、丸め、通貨、決済API、キャンセル、在庫、監査ログが含まれるかもしれない。この全体を一度に証明しようとしない。最初は、「割引後小計に送料を足す純粋関数が仕様式と一致する」のような核だけを証明する。
\end{softwarebox}

過剰形式化を避けるために、設計文書には次のような境界を書くとよい。

\begin{itemize}
\item この Lean モデルでは、金額を \texttt{Nat} として扱う。
\item 小数、丸め、通貨単位、税率変更はモデル外とする。
\item 外部決済APIの成功・失敗は統合テストで確認する。
\item この theorem は、純粋な計算核の意味保存だけを保証する。
\end{itemize}

このように境界を書くことで、「証明済み」という言葉が過大に解釈されるのを防げる。

\section{29.4 証明可能な核を抽出する}

証明可能な核は、自然言語仕様から機械的に出てくるわけではない。人間が設計判断として切り出す必要がある。実務で使いやすい手順は次のとおりである。

\begin{enumerate}
\item 仕様文から、入力、出力、状態、操作を抜き出す。
\item 外部依存を一度外し、純粋な関数として書ける部分を探す。
\item 守りたい性質を、等式、含意、不変条件、round-trip などに言い換える。
\item モデル化の前提を明示する。
\item theorem 名を仕様項目IDと対応づける。
\end{enumerate}

ここでは、出金処理の核を例にする。現実の出金には、認証、監査ログ、DB更新、並行性、通知などがある。しかし、残高と出金額だけを見た純粋な核は、\texttt{Option Nat} を返す関数として表せる。

\begin{listing}[htbp]
\begin{minted}{lean4}
def withdrawCore (balance amount : Nat) : Option Nat :=
  if h : amount <= balance then
    some (balance - amount)
  else
    none

theorem withdrawCore_success
    (balance amount : Nat) (h : amount <= balance) :
    withdrawCore balance amount = some (balance - amount) := by
  simp [withdrawCore, h]

theorem withdrawCore_failure
    (balance amount : Nat) (h : ¬ (amount <= balance)) :
    withdrawCore balance amount = none := by
  simp [withdrawCore, h]
\end{minted}
\caption{出金処理から純粋な核を切り出す}
\label{lst:ch29-withdraw-core}
\end{listing}

このコードでは、出金可能な場合と不可能な場合を定理として分けている。\texttt{withdrawCore\_success} は、出金額が残高以下なら、結果は残高から出金額を引いた値になることを言う。\texttt{withdrawCore\_failure} は、出金額が残高を超えるなら失敗することを言う。

\begin{leanbox}
ここで重要なのは、DBもログも登場しないことである。これは欠落ではなく、意図的な抽象化である。Lean では、まず純粋な残高更新ルールだけを証明する。DB更新や監査ログは、別途統合テストや運用テストで確認する。
\end{leanbox}

\subsection{圏論的な読み方}

圏論の言葉に戻すと、証明可能な核を抽出する作業は、複雑な実務世界から、小さな型と関数の世界へ写す作業である。対象は \texttt{Nat} や \texttt{Option Nat} のような型になり、射は \texttt{withdrawCore} や \texttt{chargeImpl} のような関数になる。

このとき、守りたい性質は、射の等式や保存性として書かれる。

\begin{itemize}
\item 実装が仕様と一致する：\texttt{impl x = spec x}
\item 変換して戻すと元に戻る：\texttt{rollback (migrate x) = x}
\item adapter と map が可換である：\texttt{adapter (map f x) = map f (adapter x)}
\item 操作後も不変条件を満たす：\texttt{Invariant x -> Invariant (step x)}
\end{itemize}

この形まで落とせるなら、Lean で扱う候補になる。逆に、この形に落とすために大量の外部環境を再現しなければならないなら、通常テストや監視に残す判断が自然である。

\section{29.5 設計文書に Lean theorem を埋め込む}

Lean theorem は、Lean ファイルの中だけに閉じ込めると、チームの設計判断に接続しにくい。設計文書には、自然言語の仕様項目と theorem の対応を書いておくとよい。

\begin{softwarebox}
おすすめの書き方は、「仕様項目ID」「Lean theorem」「前提条件」「保証すること」「保証しないこと」「対応するテスト」を並べる形式である。これにより、Lean の証明がどこまでの意味を持つのかをレビューで確認しやすくなる。
\end{softwarebox}

たとえば、料金計算の核に対して、次のような設計文書を置ける。

\begin{listing}[htbp]
\begin{minted}{text}
仕様項目: SPEC-CHARGE-001
主張: 実装モデル chargeImpl は仕様モデル chargeSpec と一致する。
Lean theorem: Chapter29.SPEC_CHARGE_001
前提: 金額は Nat。小数、丸め、税率、外部決済APIはモデル外。
保証: 任意の PriceInput について chargeImpl i = chargeSpec i。
保証しないこと: 実DB、外部API、画面表示、性能、通貨丸め。
対応テスト: 決済API統合テスト、丸め仕様テスト、UI表示テスト。
\end{minted}
\caption{設計文書中の theorem 対応表の例}
\label{lst:ch29-design-doc-fragment}
\end{listing}

この文書断片に対応する theorem を、Lean 側では次のように書ける。

\begin{listing}[htbp]
\begin{minted}{lean4}
/--
SPEC-CHARGE-001:
実装モデル chargeImpl は仕様モデル chargeSpec と一致する。
前提: 金額は Nat。丸め、税率、外部決済APIはモデル外。
-/
theorem SPEC_CHARGE_001 (i : PriceInput) :
    chargeImpl i = chargeSpec i := by
  exact charge_impl_matches_spec i

end Chapter29
\end{minted}
\caption{設計文書IDに対応する theorem}
\label{lst:ch29-spec-charge-001}
\end{listing}

この対応があると、レビューで次の会話ができる。

\begin{itemize}
\item 仕様項目 \texttt{SPEC-CHARGE-001} は Lean で証明されているか。
\item theorem はどのモデルの上で成り立つか。
\item モデル外の性質は、どのテストで確認するか。
\item 実装変更時に、この theorem が壊れたら何を見直すか。
\end{itemize}

\subsection{定理名を仕様項目に寄せる}

Lean では、数学的に短い theorem 名を付けたくなることがある。しかし実務で使うなら、仕様項目IDとの対応が追いやすい名前も有効である。\texttt{SPEC\_CHARGE\_001} のような名前は数学的には少し無骨だが、設計文書、Issue、Pull Request、CIログから検索しやすい。

一方で、すべての定理を仕様IDだけにすると、Lean コードが読みにくくなる。そこで、次のように分けるとよい。

\begin{itemize}
\item 開発中の補題には、\texttt{charge\_impl\_matches\_spec} のような意味のある名前を付ける。
\item 設計文書に公開する定理には、\texttt{SPEC\_CHARGE\_001} のような仕様ID名を付ける。
\item 公開定理は補題を使って短く証明する。
\end{itemize}

この分け方をすると、Lean ファイルは読みやすく、設計文書との対応も追いやすい。

\section{実装コードと Lean モデルの距離を明示する}

Lean で最も誤解されやすい点は、本番コードと Lean モデルの距離である。Lean に書いた関数が本番コードそのものではない場合、証明は「本番コードが正しい」ことを直接意味しない。証明が意味するのは、「Lean に書いたモデルでは、この性質が成り立つ」ということである。

そのため、設計文書には次の情報を書く。

\begin{itemize}
\item 本番コードのどの関数・モジュールに対応する Lean モデルか。
\item 本番コードから何を抽象化したか。
\item 本番コードと Lean モデルを同期する責任者またはレビュー観点は何か。
\item Lean で証明した性質を補うテストは何か。
\end{itemize}

\begin{warningbox}
Lean モデルと本番コードの差分を隠してはいけない。差分を隠すと、「証明済み」という言葉だけが独り歩きする。差分を明示すれば、証明は過剰な保証ではなく、設計判断を支える正確な証拠になる。
\end{warningbox}

\section{よくある誤解}

\subsection{誤解1：証明があればテストはいらない}

証明は、モデル化した範囲では強い。しかし、外部API、DB、UI、性能、運用は、モデル外であることが多い。そこはテストや監視が必要である。

\subsection{誤解2：本番コードを全部 Lean に書くべきである}

本章の方針では、まず証明可能な核を切り出す。完全検証を目指す研究・開発プロジェクトでは別の判断もあり得るが、通常のチーム導入では、小さく始めるほうが継続しやすい。

\subsection{誤解3：証明できない仕様は悪い仕様である}

証明できない仕様にも価値はある。ユーザー体験、運用判断、性能目標、法務・業務上の例外などは、自然言語、テスト、監視、レビューで扱うほうがよい場合が多い。重要なのは、証明できないことを隠さず、確認方法を分けることである。

\section{本章のまとめ}

Lean で証明する価値が高いのは、小さく、純粋で、すべての入力について保証したい性質である。

通常のテストは、外部環境、UI、性能、統合挙動、運用上の振る舞いを確認するために必要である。

過剰形式化を避けるには、証明する核、テストする周辺、文書化する前提を分ける。

自然言語仕様から、型、関数、述語、theorem に落とせる部分を抽出すると、Lean で扱いやすくなる。

設計文書には、Lean theorem、前提条件、保証すること、保証しないこと、対応テストを対応づけて残すとよい。
